\section{Conclusion}

We presented ReAct SQL, a tool-augmented agent framework that achieves state-of-the-art performance on Text-to-SQL through three key innovations: (1) a Rich Context mechanism capturing comprehensive database metadata including indexes, JOIN paths, and dialect features; (2) a unified ReAct paradigm integrating database onboarding and SQL generation; and (3) multi-database compatibility through syntax-specific guidance. Our systematic dataset calibration of Spider dev with 221 corrected annotations ensures reliable evaluation.

Experimental results demonstrate significant improvements over existing methods, with ablation studies confirming the effectiveness of each component. The framework's task-driven approach enables adaptive database analysis and self-correcting SQL generation, addressing fundamental limitations of current Text-to-SQL systems.

\subsection{Limitations and Future Work}

While ReAct SQL achieves strong performance, several directions remain for future work:

\textbf{Scalability:} Current Rich Context construction may be time-consuming for very large databases (1000+ tables). Future work could explore incremental context building or context compression techniques.

\textbf{Dynamic Schemas:} The framework assumes relatively stable schemas. Handling frequently changing schemas would require automatic context invalidation and refresh mechanisms.

\textbf{Complex Ambiguities:} While our dataset calibration addresses many ambiguities, some queries remain inherently ambiguous. Future work could integrate interactive clarification mechanisms.

\textbf{Cost Optimization:} The ReAct loop involves multiple LLM calls. Exploring more efficient reasoning strategies or smaller specialized models could reduce computational costs.

Despite these limitations, ReAct SQL represents a significant step toward reliable, high-performance Text-to-SQL systems with strong cross-database compatibility and comprehensive schema understanding.
